\documentclass[11pt,a4paper]{article}
% Drake_preamble.tex - LaTeX Preamble Template
% 作者: 謝慕揚, MD, PhD, FESC
% 
% 使用說明:
% 1. 表格請使用 longtable，不要有垂直分隔線 |
% 2. 表格內換行使用 \newline，不要使用 \\
% 3. 藥物名稱、檢驗、檢查請維持英文
% 4. Reference 格式請使用 \href 加上驗證過的 DOI

% Including essential packages for Chinese typesetting
\usepackage{xeCJK}
\usepackage{fontspec}
% Configuring Chinese font with Noto Serif CJK TC, fallback to SimSun
\IfFontExistsTF{Noto Serif CJK TC}{
  \setCJKmainfont{Noto Serif CJK TC}
  \setCJKsansfont{Noto Serif CJK TC}
  \setCJKmonofont{Noto Serif CJK TC}
}{\setCJKmainfont{SimSun}
  \setCJKsansfont{SimSun}
  \setCJKmonofont{SimSun}
}

% Including other necessary packages
\usepackage{geometry}
\usepackage{titlesec}
\usepackage{enumitem}
\usepackage{xcolor}
\usepackage{colortbl}
\usepackage{tcolorbox}
\usepackage{booktabs}
\usepackage{longtable}
\usepackage{array}
\usepackage{graphicx}
\usepackage{tikz}
\usepackage{fancyhdr}
\usepackage{multicol}
\usepackage{datetime2}
\usepackage{hyperref}
\usepackage{mdframed}
\usepackage{amsmath}
\usepackage{amssymb}
\usepackage{multirow}

\hypersetup{
    colorlinks=true,
    linkcolor=emergency_blue,
    citecolor=emergency_blue,
    urlcolor=emergency_blue,
    pdfborder={0 0 0}
}

% Defining page geometry
\geometry{
  top=2.5cm,
  bottom=2.5cm,
  left=2cm,
  right=2cm,
  headheight=25pt
}

% Adjusting topmargin to compensate for increased headheight
\addtolength{\topmargin}{-10pt}

% Defining colors
\definecolor{emergency_red}{RGB}{186,24,27}
\definecolor{emergency_blue}{RGB}{0,114,188}
\definecolor{emergency_gray}{RGB}{120,120,120}
\definecolor{highlight_yellow}{RGB}{255,250,205}

% Setting title formats
\titleformat{\section}[block]{\Large\bfseries\color{emergency_red}}{\thesection}{10pt}{}
\titleformat{\subsection}[block]{\large\bfseries\color{emergency_blue}}{\thesubsection}{8pt}{}
\titleformat{\subsubsection}[block]{\normalsize\bfseries\color{emergency_gray}}{\thesubsubsection}{6pt}{}

% Defining custom environment for important notes
\newmdenv[
    linecolor=emergency_red,
    backgroundcolor=highlight_yellow,
    linewidth=2pt,
    topline=true,
    bottomline=true,
    leftline=true,
    rightline=true,
    innertopmargin=10pt,
    innerbottommargin=10pt,
    innerleftmargin=10pt,
    innerrightmargin=10pt
]{importantbox}

% Defining note command with tcolorbox
\newcommand{\note}[1]{
    \par
    \noindent
    \begin{tcolorbox}[
        colback=emergency_blue!5!white,
        colframe=emergency_blue,
        title=\textbf{重點提醒},
        fonttitle=\bfseries,
        boxrule=0.5pt,
        arc=3pt,
        left=6pt,
        right=6pt,
        top=6pt,
        bottom=6pt
    ]
    #1
    \end{tcolorbox}
}

% Key findings box
\newcommand{\keyfinding}[1]{
    \par
    \noindent
    \begin{tcolorbox}[
        colback=yellow!10!white,
        colframe=orange!75!black,
        title=\textbf{核心發現摘要},
        fonttitle=\bfseries,
        boxrule=0.5pt,
        arc=3pt
    ]
    #1
    \end{tcolorbox}
}

% Defining custom date format for ISO 8601
\DTMnewdatestyle{iso}{
  \renewcommand{\DTMdisplaydate}[4]{\number##1-\number##2-\number##3}
  \renewcommand{\DTMDisplaydate}[4]{\number##1-\number##2-\number##3}
}
\DTMsetdatestyle{iso}
\newcommand{\mydate}{\DTMtoday}


% Custom header for this document
\pagestyle{fancy}
\fancyhf{}
\fancyhead[L]{\textcolor{emergency_red}{Circulation 教學講義}}
\fancyhead[R]{\textcolor{emergency_blue}{STORM-PE: 機械血栓抽吸術治療肺栓塞}}
\fancyfoot[C]{\textcolor{emergency_gray}{\thepage}}
\renewcommand{\headrulewidth}{1pt}
\renewcommand{\headrule}{\hbox to\headwidth{\color{emergency_red}\leaders\hrule height \headrulewidth\hfill}}

\begin{document}

\begin{center}
{\LARGE\bfseries\textcolor{emergency_red}{Circulation}}\\[0.5cm]
{\Large\textbf{Randomized Controlled Trial of Mechanical Thrombectomy With Anticoagulation Versus Anticoagulation Alone for Acute Intermediate-High Risk Pulmonary Embolism: Primary Outcomes From the STORM-PE Trial}}\\[0.3cm]
{\large 機械血栓抽吸術合併抗凝血治療 vs 單獨抗凝血治療\\急性中高風險肺栓塞：STORM-PE 試驗主要結果}\\[0.5cm]
{\normalsize Circulation. 2025; [Published Online Ahead of Print]. DOI: 10.1161/CIRCULATIONAHA.125.077232}\\[0.3cm]
{\normalsize 整理：謝慕揚 MD, PhD, FESC \quad 日期：\mydate}
\end{center}

\vspace{0.5cm}

\keyfinding{
\begin{itemize}[leftmargin=*,itemsep=2pt]
    \item STORM-PE 為首個比較機械血栓抽吸術 (mechanical thrombectomy, MT) 合併抗凝血治療與單獨抗凝血治療的隨機對照試驗
    \item CAVT 組在 48 小時內 RV/LV ratio 下降幅度顯著優於單獨抗凝血組 (29.7\% vs 13.1\%, P $<$ 0.001)
    \item 兩組安全性相當，CAVT 組無裝置相關死亡，技術成功率達 100\%
    \item CAVT 組在功能性結果方面表現更佳：6 分鐘步行距離 (6MWT) 及功能狀態量表 (PVFS) 均有顯著改善
\end{itemize}
}

\section{研究背景 (Background)}

\subsection{肺栓塞流行病學}
\begin{itemize}[leftmargin=*,itemsep=2pt]
    \item 全球每年約每 1,000 名成人中有 1 人發生 pulmonary embolism (PE)
    \item 造成顯著的社會經濟負擔
    \item 患者功能狀態與生活品質下降
\end{itemize}

\subsection{中高風險肺栓塞 (Intermediate-High Risk PE)}
\begin{itemize}[leftmargin=*,itemsep=2pt]
    \item 定義：RV/LV ratio 升高 + cardiac biomarkers 上升
    \item 死亡風險較低風險組高出 2 倍以上
    \item 有早期臨床惡化風險
\end{itemize}

\subsection{實證空白 (Evidence Gap)}
\begin{itemize}[leftmargin=*,itemsep=2pt]
    \item Anticoagulation (AC) 為目前主要治療
    \item 上一個與 AC 比較的 RCT 已超過 10 年
    \item MT vs AC 的功能性預後不明
\end{itemize}

\note{
\textbf{STORM-PE 為首個評估 CAVT (Computer Assisted Vacuum Thrombectomy) + AC vs AC alone 的 RCT}\\
由 Penumbra 公司贊助，與 The PERT Consortium 合作進行
}

\section{研究設計 (Study Design)}

\subsection{納入條件 (Inclusion Criteria)}
\begin{itemize}[leftmargin=*,itemsep=2pt]
    \item 急性 PE 臨床症狀與徵象 ($\leq$ 14 天)
    \item CTPA 顯示至少 1 條 main 或 proximal lobar pulmonary artery 有 filling defect
    \item CTPA 上 RV/LV ratio $\geq$ 1.0
    \item Cardiac biomarkers 上升
\end{itemize}

\subsection{排除條件 (Exclusion Criteria)}
\begin{itemize}[leftmargin=*,itemsep=2pt]
    \item 血流動力學不穩定或使用 ECMO
    \item CTEPH 或 CTED 發現
    \item 原發性或轉移性腦部惡性腫瘤
    \item 預期壽命 $<$ 90 天
    \item NEWS2 $\geq$ 9
\end{itemize}

\subsection{試驗架構}

\begin{longtable}{p{4cm}p{10cm}}
\toprule
\textbf{項目} & \textbf{內容} \\
\midrule
\endfirsthead
\toprule
\textbf{項目} & \textbf{內容} \\
\midrule
\endhead
\bottomrule
\endfoot
研究設計 & Prospective, multicenter, randomized controlled trial \\
\midrule
隨機分配 & 1:1 隨機分至 CAVT + AC 組或 AC alone 組 \\
\midrule
研究人數 & 100 人 (CAVT 組 47 人，AC 組 53 人) \\
\midrule
研究中心 & 22 個國際中心 (美國、波蘭、紐西蘭、加拿大) \\
\midrule
主要終點 & 48 小時內 RV/LV ratio 變化差異 (由盲性影像核心實驗室判定) \\
\midrule
次要安全終點 & 7 天內 MAE (由 clinical events committee 判定)\newline 90 天內安全性結果 \\
\midrule
次要終點 & 90 天功能性結果與生活品質 \\
\midrule
治療時程 & 於 baseline CTPA 後 24 小時內且隨機分配後 $\leq$ 12 小時開始治療 \\
\bottomrule
\end{longtable}

\section{基線特徵 (Baseline Characteristics)}

兩組在關鍵基線指標上配對良好：

\begin{longtable}{p{5cm}p{3.5cm}p{3.5cm}p{2cm}}
\toprule
\textbf{變項} & \textbf{CAVT (N=47)} & \textbf{AC (N=53)} & \textbf{P value} \\
\midrule
\endfirsthead
\toprule
\textbf{變項} & \textbf{CAVT (N=47)} & \textbf{AC (N=53)} & \textbf{P value} \\
\midrule
\endhead
\bottomrule
\endfoot
\multicolumn{4}{l}{\textbf{人口學特徵}} \\
\midrule
年齡 (歲) & 59.5 $\pm$ 13.2 & 61.2 $\pm$ 14.2 & NS \\
女性 & 18 (38.3\%) & 28 (52.8\%) & NS \\
\midrule
\multicolumn{4}{l}{\textbf{病史}} \\
\midrule
高血壓 & 21 (44.7\%) & 35 (66.0\%) & $<$ 0.05 \\
糖尿病 & 9 (19.1\%) & 9 (17.0\%) & NS \\
DVT & 30 (63.8\%) & 32 (60.4\%) & NS \\
先前 PE & 12 (25.5\%) & 10 (18.9\%) & NS \\
\midrule
\multicolumn{4}{l}{\textbf{臨床參數}} \\
\midrule
昏厥 (Syncope) & 9 (19.1\%) & 8 (15.1\%) & NS \\
NEWS2 & 3.5 $\pm$ 1.95 & 4.1 $\pm$ 2.07 & NS \\
Heart rate (bpm) & 93.2 $\pm$ 17.36 & 98.2 $\pm$ 15.87 & NS \\
Oxygen saturation (\%) & 96.0 $\pm$ 2.59 & 95.4 $\pm$ 2.44 & NS \\
RV/LV ratio & 1.63 $\pm$ 0.36 & 1.56 $\pm$ 0.35 & 0.397 \\
RMMS & 27.3 $\pm$ 3.89 & 26.1 $\pm$ 5.51 & NS \\
\bottomrule
\end{longtable}

\section{CAVT 組手術資訊 (Procedural Data)}

\begin{longtable}{p{6cm}p{8cm}}
\toprule
\textbf{項目} & \textbf{結果} \\
\midrule
\endfirsthead
\bottomrule
\endfoot
Thrombectomy time (中位數) & 25.0 分鐘 [IQR 15.0, 41.0] \\
\midrule
Procedure time (中位數) & 56.0 分鐘 [IQR 42.0, 69.0] \\
\midrule
估計失血量 & 296.5 $\pm$ 179.4 mL \\
\midrule
Technical success & 47/47 (100\%)\newline (導管能夠到達血栓並進行抽吸) \\
\midrule
mPAP 下降 & 27.3\% (平均 8.2 mmHg) \\
\midrule
Device/Procedure-related transfusion & 0/47 (0\%) \\
\midrule
Access site complications & 0/47 (0\%) \\
\bottomrule
\end{longtable}

\section{主要結果 (Primary Outcomes)}

\subsection{主要療效終點：RV/LV Ratio 變化}

\begin{importantbox}
\textbf{CAVT 組顯示優越的 RV/LV ratio 下降}
\begin{itemize}[leftmargin=*,itemsep=2pt]
    \item \textbf{CAVT 組}：1.63 $\rightarrow$ 1.11 (相對下降 \textbf{29.7\%})
    \item \textbf{AC 組}：1.56 $\rightarrow$ 1.32 (相對下降 \textbf{13.1\%})
    \item 組間絕對下降差異：$\Delta$ 0.27 (95\% CI: 0.12, 0.43)，\textbf{P $<$ 0.001}
\end{itemize}
\end{importantbox}

\begin{longtable}{p{4cm}p{3cm}p{3cm}p{3cm}}
\toprule
\textbf{RV/LV Ratio} & \textbf{CAVT (N=46)} & \textbf{AC (N=52)} & \textbf{P value} \\
\midrule
\endfirsthead
\bottomrule
\endfoot
Baseline & 1.63 $\pm$ 0.36 & 1.56 $\pm$ 0.35 & 0.397 \\
\midrule
48 小時 & 1.11 $\pm$ 0.28 & 1.32 $\pm$ 0.31 & $<$ 0.001 \\
\midrule
絕對下降 & 0.52 $\pm$ 0.37 & 0.24 $\pm$ 0.40 & $<$ 0.001 \\
\bottomrule
\end{longtable}

\subsection{血栓負荷變化 (Thrombus Burden)}

以 Refined Modified Miller Score (RMMS) 評估：
\begin{itemize}[leftmargin=*,itemsep=2pt]
    \item \textbf{CAVT 組}：相對平均下降 \textbf{42.1\%}
    \item \textbf{AC 組}：相對平均下降 \textbf{15.6\%}
    \item CAVT 組血栓負荷下降為 AC 組的 \textbf{2.7 倍}，P $<$ 0.001
\end{itemize}

\subsection{48 小時內生理參數改善}

\begin{longtable}{p{5cm}p{3cm}p{3cm}p{2.5cm}}
\toprule
\textbf{參數} & \textbf{CAVT} & \textbf{AC} & \textbf{P value} \\
\midrule
\endfirsthead
\bottomrule
\endfoot
Heart rate (bpm) at 48h & 80.0 & 86.4 & 0.022 \\
\midrule
Tachycardia (HR$>$100) at 48h & 2.2\% & 20.0\% & 0.008 \\
\midrule
Supplemental O$_2$ (L/min) at 48h & 0.5 & 1.4 & 0.027 \\
\midrule
NEWS2 at 48h & 1.8 & 2.7 & 0.034 \\
\bottomrule
\end{longtable}

\section{安全性結果 (Safety Outcomes)}

\subsection{7 天內主要不良事件 (MAE $\leq$ 7 Days)}

\begin{longtable}{p{6cm}p{3cm}p{3cm}p{2cm}}
\toprule
\textbf{事件} & \textbf{CAVT (N=47)} & \textbf{AC (N=53)} & \textbf{P value} \\
\midrule
\endfirsthead
\bottomrule
\endfoot
Composite MAE $\leq$ 7 Days & 2 (4.3\%) & 4 (7.5\%) & 0.681 \\
\midrule
Clinical deterioration requiring rescue therapy & 1 (2.1\%) & 3 (5.7\%) & 0.620 \\
\midrule
PE-related mortality & 2 (4.3\%) & 0 (0\%) & 0.218 \\
\midrule
Symptomatic recurrent PE & 0 (0\%) & 0 (0\%) & $>$ 0.999 \\
\midrule
Major bleeding (BARC 3a-5) & 1 (2.1\%) & 1 (1.9\%) & $>$ 0.999 \\
\midrule
Device/Procedure-related MAEs & 0 (0\%) & NA & NA \\
\bottomrule
\end{longtable}

\note{
\textbf{安全性重點}：
\begin{itemize}[leftmargin=*,itemsep=2pt]
    \item 兩組 90 天安全性相當
    \item \textbf{無裝置相關死亡}
    \item 7 天後無 PE 相關死亡
    \item 90 天內 symptomatic recurrent PE：CAVT 1 (2.1\%) vs AC 1 (1.9\%)
\end{itemize}
}

\section{功能性結果 (Functional Outcomes)}

\subsection{PVFS (Post-VTE Functional Status Scale)}

PVFS 為評估日常生活功能限制的序位量表：
\begin{itemize}[leftmargin=*,itemsep=2pt]
    \item CAVT 患者恢復至 PE 發作前功能狀態的可能性為 AC 組的 \textbf{2.2 倍} (OR 2.2, P = 0.032)
    \item 從 baseline 到出院的改善：OR 2.0, P = 0.067
\end{itemize}

\subsection{6 分鐘步行測試 (6MWT)}

\begin{longtable}{p{5cm}p{3.5cm}p{3.5cm}p{2cm}}
\toprule
\textbf{6MWT 距離} & \textbf{CAVT (N=27)} & \textbf{AC (N=29)} & \textbf{P value} \\
\midrule
\endfirsthead
\bottomrule
\endfoot
30 天 (公尺) & 387 & 328 & 0.099 \\
\midrule
90 天 (公尺) & 472 & 376 & \textbf{0.019} \\
\midrule
30-90 天改善量 ($\Delta$) & 85 & 48 & -- \\
\midrule
90 天預測距離達成率 & 94.0\% & 75.2\% & \textbf{0.022} \\
\bottomrule
\end{longtable}

\begin{importantbox}
\textbf{功能性結果重點}：
\begin{itemize}[leftmargin=*,itemsep=2pt]
    \item CAVT 患者在 90 天時步行距離顯著較遠
    \item 30-90 天的步行距離改善為 AC 組的 \textbf{1.8 倍}
    \item CAVT 患者能完成顯著較高比例的預測步行距離 (94\% vs 75\%)
\end{itemize}
\end{importantbox}

\subsection{呼吸困難評估}

出院時 Borg dyspnea scale 與 mMRC 呼吸困難量表：
\begin{itemize}[leftmargin=*,itemsep=2pt]
    \item Borg dyspnea score：CAVT 1.0 vs AC 1.5 (P = 0.109)
    \item mMRC $\geq$ 1：CAVT 54.5\% vs AC 73.6\% (P = 0.058)
    \item 雖未達統計顯著，但趨勢有利於 CAVT 組
\end{itemize}

\section{評估工具說明 (Assessment Tools)}

\begin{longtable}{p{3.5cm}p{10.5cm}}
\toprule
\textbf{工具} & \textbf{說明} \\
\midrule
\endfirsthead
\bottomrule
\endfoot
NEWS2 & National Early Warning Score 2：評估臨床惡化風險的綜合評分\newline 包含：respiratory rate、supplemental oxygen、oxygen saturation、systolic BP、heart rate、consciousness、temperature \\
\midrule
RMMS & Refined Modified Miller Score：評估肺動脈血栓負荷的影像學評分 \\
\midrule
RV/LV ratio & 右心室/左心室直徑比值：評估右心負荷的影像學指標 \\
\midrule
6MWT & 6-Minute Walk Test：評估功能性運動能力 \\
\midrule
PVFS & Post-VTE Functional Status Scale：評估日常生活功能限制的序位量表 \\
\midrule
Borg scale & 評估呼吸困難主觀感受的量表 (0-10 分) \\
\midrule
mMRC & Modified Medical Research Council Dyspnea Scale：評估呼吸困難對功能影響 \\
\midrule
PEmb-QoL & Pulmonary Embolism Quality of Life Questionnaire：PE 相關生活品質問卷 \\
\midrule
EQ-5D-5L & 通用健康相關生活品質問卷 \\
\bottomrule
\end{longtable}

\section{研究限制 (Limitations)}

\begin{itemize}[leftmargin=*,itemsep=2pt]
    \item 開放標籤設計 (open-label)，非盲性試驗
    \item 樣本數相對較小 (N=100)，未 powered 以偵測安全性差異
    \item 使用 surrogate endpoint (RV/LV ratio) 而非臨床 hard endpoint
    \item 短期追蹤 (90 天)
    \item 由 device 公司 (Penumbra) 贊助
\end{itemize}

\section{結論 (Conclusions)}

\begin{importantbox}
\textbf{STORM-PE 研究結論}：
\begin{enumerate}[leftmargin=*,itemsep=2pt]
    \item STORM-PE 為首個報告 MT + AC vs AC alone 結果的 RCT
    \item CAVT 組展現 \textbf{優越的 RV 擴張減少}，且安全性與 AC 組相當
    \item CAVT 治療帶來 \textbf{早期生理恢復} 及更大幅度的肺動脈阻塞減少
    \item CAVT 患者展現 \textbf{更佳的功能性改善} (PVFS、6MWT)
    \item 結果支持機械血栓抽吸術 (CAVT) 作為急性中高風險 PE 的有效治療策略
\end{enumerate}
\end{importantbox}

\section{臨床意涵 (Clinical Implications)}

\begin{itemize}[leftmargin=*,itemsep=2pt]
    \item 對於中高風險 PE 患者，CAVT 可能提供比單獨抗凝血更佳的短期結果
    \item 100\% 的技術成功率與無裝置相關輸血需求顯示良好的安全性
    \item 功能性結果改善對於患者長期生活品質具有重要意義
    \item 需要更大型試驗確認臨床 hard endpoint 的效益
\end{itemize}

\section{參考文獻 (References)}

\begin{enumerate}[leftmargin=*,itemsep=3pt]
    \item Lookstein RA, Konstantinides SV, Weinberg I, et al. Randomized Controlled Trial of Mechanical Thrombectomy With Anticoagulation Versus Anticoagulation Alone for Acute Intermediate-High Risk Pulmonary Embolism: Primary Outcomes From the STORM-PE Trial. \href{https://doi.org/10.1161/CIRCULATIONAHA.125.077232}{\textit{Circulation}. 2025; [Published Online Ahead of Print].}
    \item Rosovsky RP, Konstantinides SV, Weinberg I, et al. Primary Outcome, Functional Endpoints, and Core Lab Findings from STORM-PE. Presented at: VIVA 2025; Las Vegas, NV, USA.
    \item Enright PL, Sherrill DL. Reference equations for the six-minute walk in healthy adults. \textit{Am J Respir Crit Care Med}. 1998;158(5 Pt 1):1384-1387.
\end{enumerate}

\end{document}
